%--------------------------------------------------------------------------------
%
%
%--------------------------------------------------------------------------------
\pagebreak
\section*{WT — Window toggle}
\addcontentsline{toc}{section}{WT — Window toggle}

%--------------------------------------------------------------------------------
\begin{iPageDescEntry}{Syntax}
\texttt{wt [ <toggle> ] [ , <winNum> ]} \\
\end{iPageDescEntry}

%--------------------------------------------------------------------------------
\begin{iPageDescEntry}{Description}

The \texttt{WT} command changes the display view of a window by selecting one of
its available toggle modes.

Each window type may define one or more toggle views. The current toggle value
controls how the window content is interpreted and displayed.

When the \texttt{<toggle>} parameter is omitted, the command advances to the
next available toggle view. When the maximum toggle value is reached, the
toggle wraps around to zero.

When a toggle value is specified, the window jumps directly to that toggle
view. Toggle values range from \texttt{0} to \texttt{maxToggle}, where
\texttt{maxToggle} is window-type specific.

If the window number parameter is omitted, the command applies to the current
window.

\end{iPageDescEntry}

%--------------------------------------------------------------------------------
\begin{iPageDescOpEntry}{Example}
\begin{lstlisting}[style=iPageOpStyle]
wn mem, 0x1000      # Create a memory window.

wt                  # Advance to the next display mode.
wt                  # Advance again.

wt 0                # Jump to toggle view 0 (e.g. hex32).
wt 1                # Jump to toggle view 1 (e.g. hex64).
wt 2                # Jump to toggle view 2 (e.g. decimal).
wt 3                # Jump to toggle view 3 (e.g. ASCII).

wn dcache, 3        # Create a data cache window.
wt                  # Switch to next cache way view.
wt 1                # Display cache way 1 directly.
\end{lstlisting}
\end{iPageDescOpEntry}

%--------------------------------------------------------------------------------
\begin{iPageDescEntry}{Notes}

The meaning of toggle values is window-type specific.  
For example:
\begin{itemize}
    \item Memory windows may use toggles to switch between data formats such as
          32-bit hexadecimal, 64-bit hexadecimal, decimal, or ASCII.
    \item Cache windows may use toggles to display individual ways of a
          set-associative cache.
\end{itemize}

\end{iPageDescEntry}
